
%(BEGIN_QUESTION)
% Copyright 2014, Tony R. Kuphaldt, released under the Creative Commons Attribution License (v 1.0)
% This means you may do almost anything with this work of mine, so long as you give me proper credit

I en seriekrets kan vi formulere noen generelle prinsipper for spenning, strøm, motstand og effekt. Fullfør disse setningene, der hver setning beskriver et grunnleggende prinsipp i en seriekrets.

\vskip 10pt


\begin{itemize}
\item I en seriekrets er spenningen . . .
\item I en seriekrets er strømmen . . .
\item I en seriekrets er motstanden . . .
\item I en seriekrets er effekten . . .

\end{itemize}
\vskip 1 cm
Forklar hvorfor hver av disse reglene er sann.

\underbar{file i01140}
%(END_QUESTION)





%(BEGIN_ANSWER)

\noindent
``I en seriekrets {\it summeres spenningsfallene til totalspenningen}.''

\vskip 10pt

Dette er et uttrykk for Kirchhoffs spenningslov (KVL), der den algebraiske summen av alle spenninger i en sløyfe må være lik null.




\vskip 30pt

\noindent
``I en seriekrets {\it er strømmen lik gjennom alle komponenter}.''

\vskip 10pt

Dette er sant fordi en seriekrets per definisjon bare har én strømvei.  Siden ladningsbærere må bevege seg samlet eller ikke i det hele tatt (en konsekvens av ladningsbevaring, der elektrisk ladning verken kan skapes eller ødelegges), må strømmen målt i ett punkt i en seriekrets være den samme som strømmen målt i et annet punkt i samme krets på samme tidspunkt.



\vskip 30pt

\noindent
``I en seriekrets {\it summeres motstandene til totalmotstanden}.''

\vskip 10pt

Hver motstand i en seriekrets virker mot strømmen.  Når motstander kobles i serie, summeres motstanden til en større totalmotstand fordi den samme strømmen må gå gjennom hver motstand.



\vskip 30pt

\noindent
``I en seriekrets {\it summeres effektforbruket til totaleffekten}.''

\vskip 10pt

Dette er et uttrykk for energibevaring, som sier at energi ikke kan skapes eller ødelegges.  All effekt som forbrukes i en last i kretsen må kunne føres tilbake til kilden, uansett hvordan lastene er koblet sammen.

%(END_ANSWER)





%(BEGIN_NOTES)

Reglene for serie- og parallellkretser er svært viktige for elevene å forstå.  En tendens hos mange elever er imidlertid at de pugger reglene i stedet for å forstå dem.  Elever kan jobbe hardt med å memorere reglene uten å forstå {\it hvorfor} de er sanne, og dermed ofte mislykkes i å huske eller bruke dem riktig.

%INDEX% Elektronikkrepetisjon: serie- og parallellkretser

%(END_NOTES)

